% Chapter 1

\chapter{TYPE THE TITLE OF THE CHAPTER HERE}

\section{Introduction}

Ensure that MsWord is set and defaulted to English (UK) or English (SA). Your margins should be 3cm (left) (as this allows for binding) and 2cm (right, top and bottom). Type your dissertation/thesis for all chapters in 1½-line spacing, 11-point. (Items in the Bibliography and information typed inside tables are typed in single-line spacing.) Use \underline{one space} after all punctuation marks. Use decimal subdivisions as indicated below. Consult \textcite{van-aswegen-2006} for information on research writing and bibliographic style. The collation of the thesis is also discussed at length \parencite[33]{van-aswegen-2006}. Note the use of the Harvard style of bibliographic citation. Footnotes should be used sparingly (use Microsoft Word: Insert -- Reference -- Footnote), mainly to clarify concepts or add information which does not pertain directly to the text.\footnote{Candidates writing a thesis or dissertation will find helpful manuals in the CPUT library at DDC number 808.02.} Note the justified right-hand margin in the text. NB. Do \underline{not justify right} tables or the Bibliography/References, as this will give you unnecessary white space. Press enter twice to move to the next paragraph.

Two authors in running text: \textcite{ellis-peters-2000} argue that literature is important (Ellis and Peters, 2000:45).

Two authors in parentheses: \parencite[45]{ellis-peters-2000}.

Three or more authors: \textcite{henderson-etal-1987} claim that moral philosophy is fundamental (Henderson et al., 1987:12).

\subsection{Heading: use sentence case, as in this example}

\subsection{Heading}

\subsubsection{Heading}

\subsubsection{Heading}

% Example figure
\begin{figure}[htbp]
\centering
% \includegraphics[width=0.8\textwidth]{example-image}
\caption{Add caption \underline{below} the figure (10-point bold)}
\label{fig:1.1}
\end{figure}

\textbf{Use single-line spacing and sentence case with no full stop}

\textbf{Centre, or align left}

\textbf{(Adapted from Bloggs, 1999:34)}

% Example table
\begin{table}[htbp]
\centering
\caption{Add caption \underline{above} the table (10-point bold), single-line spacing, sentence case, and no full stop}
\begin{tabularx}{\textwidth}{|Y|Y|Y|Y|>{\raggedleft\arraybackslash}p{1.8cm}|}
\hline
\textbf{Heading} & \textbf{Heading} & \textbf{Heading} & \textbf{Heading} & \textbf{Rands} \\
\hline
Remember to justify left and use single spacing. Check your punctuation at the end and be consistent. & Use a smaller font if you have a great deal of text. This is 9-point Arial. & Eight-point is generally too small to be read comfortably. & Use right justification for figures, as in the next table. Remember to right align the caption, also. & 1 234 \\
\hline
130 & 18 987 & 120 764 & 123 & \\
\hline
Text of tables should not ``overflow'' onto the next page, but should be enclosed in a new row and column. & Very wide tables should be ``landscaped'' -- remember to use a wide-enough margin for binding. & & & \\
\hline
\end{tabularx}
\label{tab:1.1}
\end{table}

\textbf{Equations}

\begin{equation}
\label{eq:1.1}
E = mc^2
\end{equation}
See Equation~\ref{eq:1.1}.
